% =========================================================================
% PREAMBLE CẤU HÌNH BỘ ĐỀ KIỂM TRA ĐỊNH KÌ TOÁN 1 THEO TỪNG BÀI HỌC
% TRUNG TÂM: LỚP TOÁN THẦY LĨNH - WEBSITE: loptoanthaylinh.vn
% FONT CHỮ: HIỆN ĐẠI, KHÔNG CHÂN (SANS-SERIF), RÕ RÀNG CHO TIỂU HỌC
% =========================================================================

\usepackage[utf8]{vietnam}
\usepackage{amsmath,amssymb}

% Font chữ hiện đại: TeX Gyre Heros (phong cách Helvetica không chân, cực kì sắc nét)
\usepackage{tgheros}
\renewcommand*\familydefault{\sfdefault}

\usepackage{graphicx}
\usepackage{fontawesome,pgfornament}
\usepackage{enumerate}
\usepackage[unicode,hidelinks]{hyperref}
\usepackage[top=1.8cm,bottom=1.8cm,left=1.8cm,right=1.8cm]{geometry}
\usepackage{fancyhdr}
\usepackage{tabularx,booktabs,colortbl,array,multirow}

% Định nghĩa bảng màu chuẩn hiện đại
\definecolor{slate}{RGB}{100,116,139}
\definecolor{myteal}{RGB}{13,148,136}
\definecolor{darkteal}{RGB}{15,118,110}

% Gói vẽ hình TikZ
\usepackage{tikz, tkz-tab, tkz-euclide}
\usetikzlibrary{shadings, fadings, shadows, shapes, shapes.geometric, arrows, arrows.meta, calc, intersections, angles, quotes, patterns, decorations.text, decorations.pathmorphing, backgrounds, positioning, fit, petri, mindmap, matrix}

% Gói khung tcolorbox
\usepackage[most]{tcolorbox}
\tcbuselibrary{skins,breakable}

% Gói tạo mã QR Code
\usepackage{qrcode}

% Gói soạn câu hỏi trắc nghiệm chuẩn BTTeX Pro
\usepackage[dethi]{ex_test}
\renewtheorem{ex}{\faEdit~Câu}
\newtheorem{bt}{\faEdit~Bài}

% Các kí hiệu toán học chuẩn
\renewcommand{\ge}{\geqslant}
\renewcommand{\geq}{\geqslant}
\renewcommand{\le}{\leqslant}
\renewcommand{\leq}{\leqslant}
\newcommand{\hoac}[1]{\left[\begin{aligned}#1\end{aligned}\right.}
\newcommand{\heva}[1]{\left\{\begin{aligned}#1\end{aligned}\right.}

% -------------------------------------------------------------------------
% CẤU HÌNH HEADER / FOOTER TRANG IN
% -------------------------------------------------------------------------
\pagestyle{fancy}
\fancyhf{}
\fancyhead[L]{\small\color{teal!85!black}\bfseries Lớp Toán Thầy Lĩnh -- loptoanthaylinh.vn}
\fancyhead[R]{\small\color{slate}\itshape Đề kiểm tra Toán 1 theo từng bài}
\fancyfoot[C]{\bfseries\thepage}
\renewcommand{\headrulewidth}{0.6pt}
\renewcommand{\footrulewidth}{0pt}

% -------------------------------------------------------------------------
% MACRO: KHUNG TIÊU ĐỀ + HỌ TÊN + ĐIỂM + MÃ QR CHO TỪNG BÀI KIỂM TRA
% Cú pháp: \TaoPhieuKiemTra{Chủ đề}{Số bài}{Tên bài}{Link QR thi online}
% -------------------------------------------------------------------------
\newcommand{\TaoPhieuKiemTra}[4]{%
  \begin{tcolorbox}[
    enhanced,
    colback=white,
    colframe=teal!80!black,
    arc=3.5mm,
    boxrule=1.2pt,
    left=3mm, right=3mm, top=2.5mm, bottom=2.5mm,
    shadow={1.5mm}{-1.5mm}{0mm}{black!10}
  ]
    % HÀNG 1: LOGO TRUNG TÂM | TIÊU ĐỀ BÀI KIỂM TRA | MÃ QR CODE
    \begin{minipage}[c]{0.26\textwidth}
      \centering
      \begin{tcolorbox}[colback=teal!10, colframe=teal!80!black, boxrule=0.6pt, arc=1.5mm, halign=center, top=1mm, bottom=1mm, left=1mm, right=1mm]
        {\bfseries\color{teal!95!black}\small LỚP TOÁN THẦY LĨNH}\\[1pt]
        {\tiny\color{slate}\faGlobe~loptoanthaylinh.vn}\\[1pt]
        {\tiny\color{slate}\faPhone~Hotline: Thầy Lĩnh}
      \end{tcolorbox}
    \end{minipage}
    \hfill
    \begin{minipage}[c]{0.54\textwidth}
      \centering
      {\footnotesize\bfseries\color{teal!80!black}\MakeUppercase{#1}}\\[2pt]
      {\large\bfseries\color{black}\MakeUppercase{BÀI #2: #3}}\\[3pt]
      {\scriptsize\itshape Thời gian làm bài: 35 phút -- Cấu trúc 3 phần chuẩn 2025}
    \end{minipage}
    \hfill
    \begin{minipage}[c]{0.16\textwidth}
      \centering
      \qrcode[height=1.5cm]{#4}\\[2pt]
      {\tiny\bfseries\color{teal!90!black}\faQrcode~Quét QR thi}
    \end{minipage}

    \vspace{2mm}\hrule\vspace{2mm}

    % HÀNG 2: THÔNG TIN HỌC SINH & KHUNG ĐIỂM / LỜI PHÊ
    \noindent
    \begin{minipage}[t]{0.68\textwidth}
      \vspace{1mm}
      \textbf{Họ và tên học sinh:} \dotfill\\[5pt]
      \textbf{Lớp:} \dotfill\quad \textbf{Trường:} \dotfill\\[5pt]
      \textbf{Ngày kiểm tra:} \dotfill/\dotfill/202\dots
    \end{minipage}
    \hfill
    \begin{minipage}[t]{0.30\textwidth}
      \begin{tcolorbox}[
        enhanced,
        colback=teal!5,
        colframe=teal!70!black,
        boxrule=0.8pt,
        arc=2mm,
        halign=center,
        top=1.5mm, bottom=1.5mm, left=1mm, right=1mm
      ]
        {\bfseries\footnotesize ĐIỂM SỐ}\\[2pt]
        {\Large\bfseries\color{teal!90!black} \dots\dots / 10}
      \end{tcolorbox}
    \end{minipage}
  \end{tcolorbox}
  \vspace{2mm}
}

% -------------------------------------------------------------------------
% MACRO: TIÊU ĐỀ LỜI GIẢI CHI TIẾT (DÙNG TRONG PHẦN 2)
% Cú pháp: \TaoHeaderLoiGiai{Chủ đề}{Số bài}{Tên bài}
% -------------------------------------------------------------------------
\newcommand{\TaoHeaderLoiGiai}[3]{%
  \begin{tcolorbox}[
    colback=teal!10,
    colframe=teal!80!black,
    arc=2mm,
    boxrule=1pt,
    left=3mm, right=3mm, top=2mm, bottom=2mm
  ]
    \centering
    {\bfseries\color{teal!90!black}\large HƯỚNG DẪN GIẢI CHI TIẾT -- BÀI #2: #3}\\[2pt]
    {\footnotesize\color{black!80}\MakeUppercase{#1}}
  \end{tcolorbox}
  \vspace{2mm}
}

% -------------------------------------------------------------------------
% ĐỊNH NGHĨA CHẾ ĐỘ PHẦN 1 (ĐỀ BÀI) VÀ PHẦN 2 (LỜI GIẢI CHI TIẾT)
% -------------------------------------------------------------------------
\newcommand{\BatCheDoDeThi}{%
  % Ẩn hoàn toàn khối lời giải trong phần đề bài để học sinh làm bài
  \renewcommand{\loigiai}[1]{}%
}

\newcommand{\BatCheDoLoiGiai}{%
  % Hiển thị lời giải chi tiết đẹp mắt với khung và màu sắc nổi bật
  \renewcommand{\loigiai}[1]{%
    \par\vspace{1.5mm}\noindent
    \begin{tcolorbox}[
      colback=slate!5!white,
      colframe=teal!50!black,
      boxrule=0.5pt,
      arc=1.5mm,
      left=2.5mm, right=2.5mm, top=1.5mm, bottom=1.5mm
    ]
      {\small\bfseries\color{teal!85!black}\faLightbulbO~Lời giải chi tiết:}\par
      {\small ##1}
    \end{tcolorbox}
    \par\vspace{1.5mm}
  }%
}
